\section{Open questions}\label{sec:discussion}
Our bounds determine the worst-case circuit compilation rate:
$2^{\Theta(m)}$ over growing binary signatures and
$2^{\Theta(m/\log(m+2))}$ over every fixed nonempty binary signature,
where $m$ is the source syntax-tree size. Several questions remain.

\paragraph{Terms over a fixed signature.}
Theorem~\ref{thm:fixed-vocabulary} determines the optimal rate for
circuits, but its upper bound uses sharing. Does every \FOthree\
formula of size $m$ over a single binary relation admit an equivalent
\CoR\ term of size $2^{o(m)}$, or can some family require
$2^{\Omega(m)}$ composition occurrences?
Our current lower bound for terms is
$2^{\Omega(m/\log(m+2))}$, and the classical general upper bound is
$2^{O(m)}$. This question does not affect the model-checking
application, which evaluates circuits directly.

\paragraph{Domain size.}
Propositions~\ref{prop:domain-size} and~\ref{prop:small-domains}
place the onset of the blowup between
domain sizes $2^k/(8k)$ and $O(k4^k)$. Our reference structures cannot
be much smaller: for fixed distinct $u,v$, the $4^k$ pairs of colors have
disjoint sets of at least two witnesses each, so a reference structure
has at least $2\cdot4^k+2$ vertices. The upper bound, on the other hand,
uses only row projections. Locating the onset may require a different
reference structure or circuits whose compositions combine different
rows.

\paragraph{The exact composition count for $G_k$.}
On arbitrary finite inputs the optimum lies between $2^{k-1}$ and $2^k$;
without converse it is exactly $2^k$. Proposition~\ref{prop:promise}
attains $2^{k-1}$ on the reference structure and all its folded
structures. Consequently these tests alone cannot improve the full-\CoR\
lower bound. Determining the optimum on arbitrary inputs calls for
either a more efficient general circuit or tests exposing an additional
obstruction.
